\section[부록 A. 구현·검증 기록]{구현·검증 기록}

\begin{frame}[plain,label=section-appendix-a]
  \hypertarget{appendix-a}{}
  \renewcommand{\insertsectionnumber}{A}
  \sectionpage
\end{frame}

\begin{frame}[label=current-status]{같은 물리 정보를 같은 의미로 전달한다}
  \framesubtitle{모델별 내부 구성은 달라도, 수신 모듈이 읽는 형식과 규약은 하나로 맞춘다.}
  \begin{block}{개발의 중심}
    \textbf{NBCP · 사각격자·삼각격자 하이젠버그}
    $\;\longrightarrow\;$ 공통 모델 정의
    $\;\longrightarrow\;$ 계산 방법과 물리량
  \end{block}
  \medskip
  \begin{tabularx}{\textwidth}{@{}p{2.1cm}Yp{2.55cm}@{}}
    \toprule
    상태 & 현재 내용 & 다음 단계 \\
    \midrule
    \textbf{방향 합의} & 모델·계산계·기준상태·방법의 경계, 간결한 모델 작업 공간 & — \\
    \textbf{범위 결정} & 1차 범위, \texttt{terms} 형식, g-텐서와 외부 $B$, 벤치마크·최소 ED, \texttt{spintoolkit} & 필드 규약 작성 \\
    \textbf{규약 검토} & 필드·좌표·단위·항 수 규칙, 상태 정합 조건 & 세부 규약 검토 \\
    \textbf{검증 기록 있음} & NBCP 모델 추출, 파일 배치, 0--4단계 구현(2b 포함) & 검토 항목 확인 \\
    \textbf{미구현} & 위상량(5단계), 2-마그논 연속체, ED flux, TN, 여러 차원 영공간 & 5단계부터 \\
    \bottomrule
  \end{tabularx}
  \slidesource{설계 상태와 적용 범위}{contract,extraction,layout}
\end{frame}

\begin{frame}[label=implemented]{첫 구현 단계: NBCP 모델 구성 추출}
  \framesubtitle{2026-09-23 기록 · 원본 예제의 계산 동작을 유지하며 소유 위치를 분리했다.}
  \begin{columns}[T,onlytextwidth]
    \begin{column}{0.49\textwidth}
      \begin{block}{적용된 변경}
        \begin{itemize}
          \item \texttt{exchange.py}: NN·NNN 행렬 구성
          \item \texttt{unit\_cells.py}: 1--4 MSL 생성
          \item 기존 예제 import 호환 유지
          \item 모델 전용 호출자 3곳의 import 정리
          \item 연구 기록의 소스 해시 목록 갱신
        \end{itemize}
      \end{block}
    \end{column}
    \begin{column}{0.47\textwidth}
      \begin{block}{이 단계의 범위}
        기존 7개 함수의 본문·시그니처를 보존했다.\\[0.8em]
        공통 \texttt{SpinModel} 자료형, 후보 상태 분리,
        결과 재사용은 후속 단계다.
      \end{block}
      \smallskip
      \supportingtext{작업 전 최신 원본의 코드 7개와 테스트 파일 7개를 기준선으로
      가져왔다. 이 기존 솔버 수정들은 이번 모델 추출의 신규 변경과 구분한다.}
    \end{column}
  \end{columns}
  \slidesource{기준선·호환성·함수 비교 기록}{extraction}
\end{frame}

\begin{frame}[label=verification]{검증 결과와 그 적용 범위}
  \framesubtitle{2026-09-23 첫 모델 추출 단계의 기록 · 현행 폴더 이행 검증은 다음 쪽에 기록한다.}
  \begin{tabularx}{\textwidth}{@{}p{4.0cm}p{2.7cm}Y@{}}
    \toprule
    확인 항목 & 기록된 결과 & 범위 \\
    \midrule
    모델·해밀토니안 전후 비교 & 40건, 차이 0 & 4 후보, 5 파라미터군, 고정·난수 각도 \\
    후보 전체 탐색 전후 비교 & 4회, 차이 0 & $N=3$, 명시한 각도 제약, classical·MAGSWT \\
    Legacy 대조 & 통과 & 모델 12 조합, 고전 에너지 36 평가 \\
    신규 모델 테스트 & 17 통과 & 데이터·에너지·선택 결합·import 호환 \\
    전체 테스트 & 215 통과 / 5 실패 & 변경 전 198 통과 / 동일한 5 실패 \\
    \bottomrule
  \end{tabularx}
  \medskip
  \begin{mdcallout}[warning]{판정의 경계}
    기존 5개 실패는 \texttt{CommensurateStructure}의 각도 수와 셀 크기 불일치다.
    Legacy 대조는 알려진 DM 회전·양자 해밀토니안 차이를 제외한다.\\
    이 결과는 동작 보존의 근거이며, 상도 수렴이나 물리 모델 승인을 뜻하지 않는다.
  \end{mdcallout}
  \slidesource{이전 모델 추출 단계의 시험 조건과 실패 목록}{extraction}
\end{frame}

\begin{frame}[label=layout-verification]{두 번째 구현 단계: 기능별 패키지 정리}
  \framesubtitle{2026-09-23 · 기존 함수·수치·공개 호출을 보존하며 파일 소유 위치를 옮겼다.}
  \begin{tabularx}{\textwidth}{@{}p{3.4cm}p{3.1cm}Y@{}}
    \toprule
    확인 항목 & 결과 & 확인 범위 \\
    \midrule
    구현 본문 비교 & 16개 모듈 보존 & import·문서 문자열을 제외한 구문 비교 \\
    상수·기본값 분리 & 20개 정의 보존 & 물리상수·수치 설정·기저 행렬 \\
    수치 결과 전후 대조 & 40건·4회, 차이 0 & 모델·행렬·에너지·최적화 결과 \\
    전체 테스트 & 217 통과 / 5 실패 & 변경 전 215 통과 / 동일한 5 실패 \\
    기존 호출 호환 & 2개 테스트 통과 & 옛 모듈 17개 경로·저장 객체 복원 \\
    \bottomrule
  \end{tabularx}
  \par\medskip
  \textbf{현재 남은 경계}\\[0.4em]
  \texttt{SpinSystem}은 여전히 각도를 포함하고, 에너지 평가는 고전 항과
  스핀파 보정을 함께 다룬다. 새 \texttt{SpinModel} 전달 규약은 후속 구현이다.
  \slidesource{격리한 준비 공간에서 검증한 폴더 이행 기록}{layout}
\end{frame}

\begin{frame}[label=rename-verification]{0단계: 패키지 이름 변경}
  \framesubtitle{2026-09-29 · \texttt{lswt}$\to$\texttt{spintoolkit}, \texttt{methods/spin\_wave}$\to$\texttt{methods/lswt}. 계산 코드는 바꾸지 않았다.}
  \begin{tabularx}{\textwidth}{@{}p{3.4cm}p{3.4cm}Y@{}}
    \toprule
    확인 항목 & 결과 & 확인 범위 \\
    \midrule
    전체 테스트 & 220 통과 / 5 실패 & 변경 전 217 / 같은 5 실패. 차이는 교체한 호환 테스트(1개$\to$4개)뿐 \\
    수치 결과 전후 대조 & 208개 값, 차이 0 & 40건의 $\mathsf H(\mathbf k)$·고전 에너지, 16개 탐색 결과. 스크립트를 저장소에 보존 \\
    옛 이름 호환 & 4개 테스트 통과 & \texttt{lswt.*}·옛 \texttt{core}·\texttt{solvers} 경로, 옛 pickle, DeprecationWarning \\
    배포·설치 & wheel, editable 설치 & 네트워크 없이 빌드; 저장소 밖에서 두 이름 import \\
    \bottomrule
  \end{tabularx}
  \par\medskip
  \supportingtext{기준 커밋 \texttt{49f047a}에서 비교했다. 다른 미커밋 수정이 있는 README·CLAUDE.md·AGENTS.md·정책 문서의 경로는 아직 갱신하지 않았다.
  탐색 중 \texttt{find\_minimum}이 \texttt{angle\_setting=None}을 처리하지 못하는 기존 결함을 발견했다.}
  \slidesource{패키지 이름 변경 검증 기록}{rename}
\end{frame}

\begin{frame}[fragile,label=stage1-api]{1단계: 공통 자료형의 사용 예}
  \framesubtitle{2026-09-29 · 모델은 해밀토니안, 상태·외부 조건·계산격자 조건은 별도 객체다(D16).}
  \begin{columns}[T,onlytextwidth]
    \begin{column}{0.56\textwidth}
\begin{semiverbatim}\scriptsize
import spintoolkit as stk
from spintoolkit.methods.classical import (
    classical_energy, torques)

model = stk.models.triangular_heisenberg(J=1.0, S=0.5)
state = stk.models.state_120(model)
cond  = stk.ExternalConditions(field=(0, 0, 0.0))
geom  = stk.CalculationGeometry.finite_torus(
            [[3, 0], [0, 3]])

stk.validate_spin_state(state, model, geom)
classical_energy(model, state, cond)  # -0.375
\end{semiverbatim}
    \end{column}
    \begin{column}{0.40\textwidth}
      \small
      \begin{tabular}{@{}ll@{}}
        \toprule
        파일 & 내용 \\
        \midrule
        \texttt{system/model.py} & \texttt{SpinModel}, 검증 \\
        \texttt{system/conditions.py} & $B$, $T$, 단위 \\
        \texttt{system/geometry.py} & 계산격자 조건 \\
        \texttt{states/spin\_state.py} & \texttt{SpinState} \\
        \texttt{methods/classical.py} & 에너지·토크 \\
        \texttt{models/heisenberg.py} & 벤치마크 \\
        \bottomrule
      \end{tabular}
    \end{column}
  \end{columns}
  \par\medskip
  \supportingtext{벤치마크 해밀토니안 생성 함수(\texttt{square\_heisenberg} 등)는 \texttt{SpinModel}을,
  기준 스핀 배열 생성 함수(\texttt{polarized\_state}, \texttt{neel\_state}, \texttt{state\_120})는
  해석적으로 알려진 \texttt{SpinState}를 만든다. 배열이 바닥상태인지는 조건에 따라 다르다.}
  \slidesource{1단계 구현 기록}{stage1}
\end{frame}

\begin{frame}[label=stage1-verification]{1단계: 검증 결과}
  \framesubtitle{같은 고전 에너지 함수가 모델명 분기 없이 세 모델을 읽는다. 기존 계산 코드는 바꾸지 않았다.}
  \small
  \begin{tabularx}{\textwidth}{@{}p{3.6cm}p{3.3cm}Y@{}}
    \toprule
    확인 항목 & 결과 & 확인 범위 \\
    \midrule
    정확한 고전 에너지 & 일치, 토크 0 & 사각 Néel $-2JS^2$, 삼각 120° $-\tfrac32JS^2$, 편극상 $zJS^2-hS$ \\
    국소장 & 수치 미분과 일치 & DM·비등방 $g$·비대각 초격자의 두 사이트 모델 \\
    NBCP 교차 확인 & 차이 $<10^{-15}$ & 기본 셀 fixture 대 기존 \texttt{EnergyFunction} One MSL; 장 변환, D13 변위 해석 \\
    규약 검증 사례 & 새 테스트 53개 통과 & 단위·중복·kind·endpoint·분수 셀·배열 보호·초격자·토러스 정합·주기 최소성 \\
    기존 테스트 & 225개 결과 불변 & 273 통과 / 같은 5 실패 \\
    수치 스냅샷 & 208개 값, 차이 0 & 0단계와 같은 대조 \\
    \bottomrule
  \end{tabularx}
  \par\medskip
  \supportingtext{후속: LSWT와 \texttt{SpinModel}의 연결(2단계), 결과 머리부·JSON 직렬화(4단계), 유한 토러스 항 전개(3단계).}
  \slidesource{1단계 검증 기록}{stage1}
\end{frame}

\begin{frame}[label=stage2-verification]{2단계: NBCP 연결과 D13 확정}
  \framesubtitle{2026-09-29 · 새 모델·상태를 기존 \texttt{SpinSystem}으로 바꿔 기존 LSWT를 그대로 쓴다(D21).}
  \small
  \begin{tabularx}{\textwidth}{@{}p{3.4cm}p{3.3cm}Y@{}}
    \toprule
    확인 항목 & 결과 & 확인 범위 \\
    \midrule
    $\mathsf H(\mathbf k)$ 원소 단위 & 최대 $1.7\times10^{-16}$ & 새 모델+후보 상태 대 기존 \texttt{one\_msl}…\texttt{four\_msl}; 4셀 $\times$ 5파라미터군, 장·DM·NNN \\
    음성 대조 & 차이 $>10^{-3}$ & $d$ 부호를 뒤집으면 Three MSL의 $\mathsf H(\mathbf k)$가 달라짐 \\
    왕복 변환 & 최대 $1.7\times10^{-17}$ & 기존 \texttt{SpinSystem} $\to$ 공통 형식 $\to$ \texttt{SpinSystem} \\
    게이지 무관 양 & 일치 & 표준 범위 밖 각도에서 $E_{\mathrm{cl}}$($10^{-14}$), $E_{\mathrm{qm}}$($10^{-12}$) \\
    테스트·스냅샷 & 348 통과 / 같은 5 실패; 차이 0 & 새 테스트 72개; 기존 225개 결과 불변 \\
    \bottomrule
  \end{tabularx}
  \par\medskip
  \supportingtext{파일: \texttt{model/nbcp/model.py}(\texttt{build\_model}, 파라미터 세트, \texttt{candidate\_state}), \texttt{system/conversion.py}.
  문헌 세트는 arXiv:2505.06398 Table 1($J_z=0.1225$, $J_{xy}=0.0779$ meV, $g_c=4.716$, $g_{ab}=4.200$).
  사이트 장이 다른 \texttt{SpinSystem}은 변환하지 않는다.}
  \slidesource{2단계 검증 기록}{stage2}
\end{frame}

\begin{frame}[fragile,label=stage2b-api]{2b단계: 고전 manifold 위 영점 에너지 선택}
  \framesubtitle{2026-09-29 · D17을 새 자료형 위에 구현하고 \texttt{quantum} 경로를 삭제했다(D18, D22).}
  \begin{columns}[T,onlytextwidth]
    \begin{column}{0.55\textwidth}
\begin{semiverbatim}\scriptsize
from spintoolkit.methods import state_selection as sel

result = sel.select_on_manifold(
    model, state, conditions,
    sel.lswt_zero_point_energy("Hex_30", N=6),
    criteria=sel.SelectionCriteria(mode="physics"))

result.verdict      # "selected"
result.axis, result.phi, result.state
result.diagnostics  # 판정에 쓴 값 전부
\end{semiverbatim}
    \end{column}
    \begin{column}{0.42\textwidth}
      \small
      \begin{enumerate}\setlength{\itemsep}{0.1em}
        \item 고전 정밀화: 접평면 L-BFGS-B, 평평한 방향을 뺀 Newton
        \item 해석적 Hessian, 회전 생성자 SVD로 축 $\hat n$
        \item 궤도 36점, 12차 Fourier fit, fit 급수의 최소
        \item $E_{\mathrm{cl}}$ 일정성, $C_{\mathrm{cl}}$ 대 $C_{\mathrm{qm}}$
      \end{enumerate}
    \end{column}
  \end{columns}
  \par\medskip
  \small 판정: \texttt{selected} · \texttt{no\_selection} · \texttt{unresolved} · \texttt{competition} · \texttt{no\_degeneracy} · \texttt{not\_flat} · \texttt{axis\_required} · \texttt{unsupported\_manifold}
  \par\medskip
  \supportingtext{Hessian: $\partial^2E/\partial x_{ia}\partial x_{jb}=S_iS_j\,\mathbf e_{ia}\!\cdot\!K_{ij}\mathbf e_{jb}+\delta_{ij}\delta_{ab}S_i\,\mathbf n_i\!\cdot\!\mathbf h_i$; 중앙 차분과 $7.8\times10^{-10}$(차분 오차 수준) 안에서 일치.
  $E_{\mathrm{qm}}$은 호출 가능 객체로 받고, 기본 제공자는 변환을 거친 기존 \texttt{EnergyFunction}이다(native는 4단계).
  \texttt{quantum}을 요청하면 \texttt{select\_on\_manifold}를 안내하는 \texttt{ValueError}.}
  \slidesource{2b단계 검증 기록}{stage2b}
\end{frame}

\begin{frame}[label=stage2b-coefficients]{2b단계: 물리 근거 모드의 계수(검토 요청)}
  \framesubtitle{D19가 구현 시 정하도록 남긴 값이다. 기본값은 \texttt{definitions/defaults.py}, 실행별 값은 \texttt{SelectionCriteria}.}
  \small\renewcommand{\arraystretch}{1.1}
  \begin{tabularx}{\textwidth}{@{}P{2.6cm}YP{1.6cm}@{}}
    \toprule
    항목 & 물리 근거 모드 & 기본값 \\
    \midrule
    상태 정확도 $\delta$ & 평평하지 않은 Hessian 방향에 남은 Newton 보정의 최대 성분 & --- \\
    평평함의 곡률 하한 & $a\cdot\max|w|\cdot(\delta+r\epsilon)$ & $a=10$, $r=10^3$ \\
    생성자 rank·대칭 허용오차 & $a\cdot(\delta+r\epsilon)$; 대칭은 $[K_{\hat n},J]=0$, $g^{\mathsf T}\mathbf b\parallel\hat n$ & 같음 \\
    해상 기준 & 최대 조화 진폭 $>\max(10\times\text{fit 잔차},\,r\epsilon\max|E_{\mathrm{qm}}|)$ & 10, $10^3$ \\
    고전·양자 비교 & Hessian이 평평하면 궤도 변동비 $\Delta E_{\mathrm{cl}}/\Delta E_{\mathrm{qm}}$, 아니면 $C_{\mathrm{cl}}/C_{\mathrm{qm}}$; $[0.1,10]$이면 경쟁 & $(0.1,10)$ \\
    고전 고정 선별 & 평평하지 않은 방향: 궤도 한 칸의 고전 변화가 $E_{\mathrm{qm}}$ 변화의 10배를 넘으면 궤도 계산 없이 \texttt{no\_degeneracy} & 대역 상한 \\
    \bottomrule
  \end{tabularx}
  \par\medskip
  \supportingtext{고정 상수 모드는 간격비 $10^{-3}$, rank $10^{-4}$, 진폭 $\max(10\times\text{잔차},10^3\epsilon|E_{\mathrm{qm}}|)$ 그대로다. 두 모드 모두 궤도 36점·12차 fit이며, 판정에 쓴 값을 \texttt{diagnostics}에 기록한다.}
  \slidesource{2b단계 검증 기록}{stage2b}
\end{frame}

\begin{frame}[label=stage2b-verification]{2b단계: NBCP Y·V 상태에서 선택을 재현한다}
  \framesubtitle{Y 0.2 T, V 1.4 T, $J_{\mathrm{PD}}$ 또는 $J_\Gamma=0.010$ meV, $10^{-3}$ rad 떨어진 출발. 기준: 독립 스캔 $N=48$.}
  \small
  \begin{tabularx}{\textwidth}{@{}YP{1.1cm}P{2.0cm}P{2.0cm}P{1.6cm}@{}}
    \toprule
    상태·결합 & 조화 & 진폭/기준 ($N=6$, 12) & $C_{\mathrm{qm}}$/기준 & 해상 여유 \\
    \midrule
    Y, $J_{\mathrm{PD}}$ & 6 & 1.011, 1.002 & 1.009, 1.002 & $>3\times10^{4}$ \\
    Y, $J_\Gamma$ & 6 & 1.000, 1.000 & 1.000, 1.000 & $9.0\times10^{2}$ \\
    V, $J_{\mathrm{PD}}$ & 6 & 1.006, 1.001 & 1.005, 1.000 & $>1.6\times10^{3}$ \\
    V, $J_\Gamma$ & 3 & 1.000, 1.000 & 1.000, 1.000 & $>3.6\times10^{8}$ \\
    \bottomrule
  \end{tabularx}
  \par\medskip
  16회(두 모드) 모두 \texttt{selected}, 축 $\hat z$. 정밀화 뒤 토크 $10^{-17}$, 영 고윳값 $10^{-18}$(다음 $5\times10^{-3}$), 축 오차 $\le10^{-15}$(유한차분 시제품 $10^{-8}$).
  \par\medskip
  Y의 $J_\Gamma$: 진폭 $1.64\times10^{-12}$, fit 잔차 $1.1\times10^{-16}$, 해상 기준 $1.8\times10^{-15}$. 최소 위치는 기준과 $5.3\times10^{-6}$ rad 이내(fit 위치 분해능 $\approx10^{-5}$); 다른 경우는 $6\times10^{-7}$ 이내.
  \par\medskip
  \supportingtext{MAGSWT 정규화 값은 모든 궤도 점에서 하한 $10^{-9}$에 머물렀다. 테스트 373 통과 / 같은 5 실패(새 29개, 삭제한 \texttt{quantum} 테스트 8개), 208개 스냅샷 차이 0.}
  \slidesource{\texttt{examples/nbcp\_state\_selection\_check.py}}{stage2b}
\end{frame}

\begin{frame}[label=stage2b-controls]{2b단계: 대조군·벤치마크·MAGSWT 재현}
  \framesubtitle{DE 출발 세 seed $\times$ 두 모드. 판정이 다른 경우는 장을 기울인 대조군 하나다.}
  \small
  \begin{tabularx}{\textwidth}{@{}YP{2.9cm}P{2.9cm}@{}}
    \toprule
    경우 & 물리 근거 모드 & 고정 상수 모드 \\
    \midrule
    정확한 U(1) & \texttt{no\_selection}(대칭) & \texttt{no\_selection} \\
    $J_{\mathrm{PD}}$, $J_\Gamma$, 회전한 $J_{\mathrm{PD}}$ & \texttt{selected}, 축 오차 $\le10^{-15}$ & 같음 \\
    장을 기울임 $(0.3h,0,h)$ & \texttt{competition}, $C_{\mathrm{cl}}/C_{\mathrm{qm}}=0.46$ & \texttt{no\_degeneracy}(간격비 $2\times10^{-2}$) \\
    편극 $20h$; 정사각 편극 $h=5>h_{\mathrm{sat}}$ & \texttt{no\_degeneracy}, rank 2 & 같음 \\
    삼각 120°, $h=0$(평평한 회전 3개) & \texttt{axis\_required}; $\hat z$를 주면 \texttt{no\_selection} & 같음 \\
    \bottomrule
  \end{tabularx}
  \par\smallskip
  \textbf{MAGSWT 재현(D18)}: 출발 azimuth가 $k\pi/6$ 격자 위이면 네 경우 모두 같은 총에너지(차이 $\le4.7\times10^{-16}$),
  격자 밖(0.3 rad)이면 궤도 탐색이 $1.3\times10^{-12}$--$8.1\times10^{-6}$ meV 낮다.
  \par\smallskip
  \supportingtext{후속: 기울인 장의 \texttt{competition}은 곡률 비 규칙의 결과다. 이 규칙의 대체안(D28 검토안)과 강체 궤도의 과대평가는 부록 B.}
  \slidesource{2b단계 검증 기록}{stage2b}
\end{frame}

\begin{frame}[label=stage3-torus-rules]{3단계: 유한 토러스 전개는 모든 방법에 같은 규칙을 쓴다}
  \framesubtitle{2026-09-30 · \texttt{system/cluster.py}의 \texttt{expand\_on\_torus}가 한 번 전개하고, 고전·ED·LSWT 비교가 그 결과를 쓴다(D23).}
  \small
  \begin{tabularx}{\textwidth}{@{}P{2.9cm}YP{2.6cm}@{}}
    \toprule
    경우 & 규칙 & 근거 \\
    \midrule
    겹친 결합(2$\times$2 정사각의 $\pm x$ 이웃) & 셀마다 항을 두고 모두 합산 & 주기 클러스터의 해밀토니안; 토러스 운동량의 $\mathsf H(\mathbf k)$와 같음 \\
    한 사이트로 접히는 결합(1$\times$1, 1$\times$3 등) & 고전·ED·LSWT 모두 거부(\texttt{ClusterError}) & $S_i\!\cdot\!J\!\cdot\!S_i$는 온사이트 항; \texttt{bilinear}로 일관 처리 불가 \\
    토러스 운동량 & $\mathbf k=2\pi A^{-1}\mathbf q$, $L\mathbf q\in\mathbb Z^2$ & 셀 수만큼의 운동량 \\
    ED 운동량 부호 & $T_R|\psi\rangle=e^{-\mathrm i\mathbf k\cdot\mathbf R}|\psi\rangle$ & $a_{\mathbf k}^\dagger|0\rangle$가 $+\mathbf k$(D13) \\
    \bottomrule
  \end{tabularx}
  \par\medskip
  \texttt{classical\_energy(model, state, conditions, geometry)}: 토러스를 주면 같은 전개로 계산한다. 정합 상태에서는 열역학 극한 값과 같다(차이 0).
  \par\medskip
  \supportingtext{온사이트 이차항 kind를 도입하면 접히는 결합을 세 방법에서 그 kind로 처리할지 함께 정한다(미결 사항).}
  \slidesource{3단계 검증 기록}{stage3}
\end{frame}

\begin{frame}[fragile,label=stage3-ed-design]{3단계: ED는 일반 해밀토니안이 기본이고 대칭은 선택이다}
  \framesubtitle{키타에프·SOC NBCP처럼 U(1)이 없는 모델도 같은 경로로 푼다. 1-마그논 비교는 그 한 사용 예다(D10, D23).}
  \begin{columns}[T,onlytextwidth]
    \begin{column}{0.52\textwidth}
\begin{semiverbatim}\scriptsize
from spintoolkit.methods.ed import solve_ed, EDSector

r = solve_ed(model, geometry, conditions,
             EDSector(axis=(0, 0, 1),
                      magnon_number=1,
                      momenta="all"))
r.blocks[0].energies     # 토러스 전체 에너지
r.excitations(r.blocks[0])  # E - E_ref
r.per_site()
\end{semiverbatim}
    \end{column}
    \begin{column}{0.45\textwidth}
      \small
      \begin{tabular}{@{}P{1.5cm}P{3.6cm}@{}}
        \toprule
        층 & 내용 \\
        \midrule
        기저 & 사이트당 $2S+1$, 임의 $S$ \\
        해밀토니안 & $S^\pm,S^z$ 사다리 계수로 sparse 구성 \\
        섹터(선택) & U(1) 자화(n-마그논), 병진 운동량 \\
        풀이 & dense 또는 Lanczos \\
        \bottomrule
      \end{tabular}
    \end{column}
  \end{columns}
  \par\medskip
  \small n-마그논 섹터: 축 $\hat n$ 방향 총자화 $M=NS-n$인 블록. 축이 $\hat z$가 아니면 문제 전체를 전역 회전한다. 섹터를 바꾸는 항이 있으면 \texttt{SectorError}.
  \par\medskip
  \supportingtext{결과 본문 \texttt{EDResult}: 섹터, 블록별 운동량·차원·고유값·풀이 방법·잔차, 선택적 고유벡터, 편극 기준 에너지, 진단. 공통 머리부·JSON은 4단계.
  2-마그논 섹터, 키타에프 flux 섹터·정확해 대조는 후속.}
  \slidesource{3단계 검증 기록}{stage3}
\end{frame}

\begin{frame}[label=stage3-verification]{3단계: 검증 결과}
  \framesubtitle{기존 계산 코드는 바꾸지 않았다. LSWT는 변환을 거친 기존 $\mathsf H(\mathbf k)$를 정규화 없이 대각화했다.}
  \small
  \begin{tabularx}{\textwidth}{@{}P{3.0cm}P{2.9cm}Y@{}}
    \toprule
    확인 항목 & 결과 & 확인 범위 \\
    \midrule
    1-마그논 ED = LSWT & 최대 $1.6\times10^{-14}$ & 9경우 $\times$ 모든 토러스 운동량: 정사각·삼각 $S=\tfrac12,1,\tfrac32$, DM, 기울인 축, 벌집격자, NBCP XXZ \\
    운동량 부호 & $-\mathbf k$와 짝지으면 1.20 & DM이 있는 경우 \\
    편극 상태 & 고전 에너지와 $7\times10^{-15}$ & 갭 $h-h_{\mathrm{sat}}$ 재현 \\
    대칭 없는 해밀토니안 & 독립 구현과 $4.6\times10^{-14}$ & 무작위 $J$, NBCP+$J_{\mathrm{PD}}$+$J_\Gamma$ \\
    4$\times$4 정사각 $S=\tfrac12$ & $E_0/N=-0.7017802$ & 알려진 16-사이트 값 \\
    테스트·스냅샷 & 404 통과 / 같은 5 실패 & 새 31개; 스냅샷 차이 0 \\
    \bottomrule
  \end{tabularx}
  \par\medskip
  \supportingtext{U(1)이 깨진 NBCP($J_\Gamma$)에 1-마그논 섹터를 요청하면 \texttt{SectorError}. 실행 환경 메모(OpenBLAS 스레드)는 검증 기록에 있다.}
  \slidesource{\texttt{examples/benchmark\_stage3\_check.py}}{stage3}
\end{frame}

\begin{frame}[fragile,label=stage4a-api]{4a: 새 모델 위의 LSWT와 공통 결과 머리부}
  \framesubtitle{2026-09-30 · \texttt{solve\_lswt}는 변환을 거쳐 기존 $\mathsf H(\mathbf k)$를 쓰고, 기존 \texttt{LSWTSolver}는 그대로 둔다(D24).}
  \begin{columns}[T,onlytextwidth]
    \begin{column}{0.52\textwidth}
\begin{semiverbatim}\scriptsize
from spintoolkit.methods.lswt import (
    solve_lswt, LSWTSettings)

r = solve_lswt(model, state, conditions,
               settings=LSWTSettings(mesh=(24, 24)))
r.ground_state_energy, r.ordered_moments()
r.bands()                  # (nk, Ns)
r.hamiltonians, r.eigenvalues, r.eigenvectors
r.header.diagnostics       # 정상성, 영모드, 정규화
save_json(r, "lswt.json")  # 머리부 + 요약
\end{semiverbatim}
    \end{column}
    \begin{column}{0.45\textwidth}
      \small
      \begin{tabular}{@{}P{1.6cm}P{3.4cm}@{}}
        \toprule
        항목 & 규칙 \\
        \midrule
        k 격자 & 자기 역격자 셀, 반 칸 이동(Γ 회피); 토러스·명시적 k \\
        정규화 & 기본 없음; 영모드·음의 모드는 오류; MAGSWT는 선택 \\
        보관 & $\mathsf H(\mathbf k)$, Colpa 고유값, paraunitary $T$ \\
        JSON & 머리부와 요약; 배열은 요청 시 \\
        \bottomrule
      \end{tabular}
    \end{column}
  \end{columns}
  \par\medskip
  \supportingtext{대각화 정보는 같은 계산 안에서 4b 열역학·4c 상관함수가 재사용한다.
  발견: 정확한 Goldstone 영모드가 반올림으로 $\min\mathrm{eig}\,\mathsf H=+5.6\times10^{-16}$이 되면 Cholesky가 통과해 $\omega\sim4\times10^{-8}$, $\langle n\rangle\sim10^6$이 나왔다. 최대 고윳값의 $10^{-10}$배 이하를 영모드로 판정해 막는다.}
  \slidesource{4a 검증 기록}{stage4a}
\end{frame}

\begin{frame}[label=stage4a-verification]{4a: 검증 결과}
  \framesubtitle{해석적 LSWT 값과는 정확히, QMC·DMRG와는 근사 비교로 구분한다(D23).}
  \footnotesize\renewcommand{\arraystretch}{1.08}
  \begin{tabularx}{\textwidth}{@{}P{3.1cm}P{3.3cm}Y@{}}
    \toprule
    확인 항목 & 결과 & 기준 \\
    \midrule
    기존 \texttt{LSWTSolver} 재현 & 에너지 $10^{-17}$, 밴드 0, $\langle n\rangle$ $10^{-14}$ & NBCP 1--4 MSL, XXZ·SOC, 같은 k·MAGSWT \\
    정사각 Néel $S=\tfrac12$ & $E=-0.6579473$, $\Delta S=0.196602$ & 독립 구적; $N=128$, $\Delta S$는 $1/N$ 외삽 \\
    삼각 120° $S=\tfrac12,1$ & $E=-0.5388090$, $\Delta S=0.2613033$ & Chernyshev–Zhitomirsky: 0.436824, 0.2613032 \\
    QMC(정사각) & 에너지 1.7\%, 모멘트 1.3\% 차이 & Sandvik: $-0.669442$, $0.307447$ \\
    DMRG(삼각 모멘트) & LSWT 0.2387가 16\% 큼 & White–Chernyshev: $0.205(15)$ \\
    NBCP Y·V & Γ 격자는 영모드 보고; 이동 격자는 정규화 없이 & LSWT 차수의 우연 영모드 \\
    테스트·스냅샷 & 421 통과 / 같은 5 실패 & 새 17개; 스냅샷 차이 0 \\
    \bottomrule
  \end{tabularx}
  \slidesource{\texttt{examples/lswt\_stage4a\_check.py}}{stage4a}
\end{frame}

\begin{frame}[label=stage4b-zero-modes]{4b: 영모드는 세 구간으로 판정하고, 애매하면 사용자가 정한다}
  \framesubtitle{2026-09-30 · $\omega$가 아니라 $\lambda_{\min}(\mathsf H(\mathbf k))/\text{scale}$로 판정한다. 작은 갭 $\Delta$는 $\lambda\sim\Delta^2/\text{대역폭}$이다(D25).}
  \footnotesize\renewcommand{\arraystretch}{1.08}
  \begin{tabularx}{\textwidth}{@{}P{3.3cm}P{3.6cm}Y@{}}
    \toprule
    경우 & 판정 & 비고 \\
    \midrule
    정사각 Néel, 삼각 120° & 영(Γ), Goldstone & 깨진 대칭 2개, 3개 \\
    이방성 $10^{-2}$, $10^{-5}$ & 갭 0.2835, 0.00894 & 해석값 $2\sqrt{(1+d)^2-1}$과 일치 \\
    이방성 $10^{-7}$, $10^{-9}$ & \textbf{후보}(갭 추정 $9\times10^{-4}$, $9\times10^{-5}$) & 사용자가 \texttt{gapless}를 정해야 진행 \\
    J1--J2, $J_2=J_1/2$ & 영, 선(12방향 중 1방향 평평) & Γ 밖; 이동 격자로도 피할 수 없음 \\
    NBCP Y: XXZ / $J_{\mathrm{PD}}$, $J_\Gamma$ & 영: Goldstone / \textbf{우연} & 대칭 없음, 고전 평평 방향 1개 \\
    편극 상태 & 갭 $h-h_{\mathrm{sat}}$ & 후보 없음 \\
    \bottomrule
  \end{tabularx}
  \par\medskip
  \supportingtext{구간: 영 $\le10^{-12}$, 후보 $\le10^{-6}$, 그 위는 갭. 탐색점: 자기·원시 격자의 고대칭점, 격자 최소점에서 연속 최소화, 12방향 선 검사.
  후보가 있으면 \texttt{thermal\_quantities}는 \texttt{ZeroModeCandidateError}로 멈추고, 사용자의 결정은 결과 머리부에 기록된다.}
  \slidesource{4b 검증 기록}{stage4b}
\end{frame}

\begin{frame}[label=stage4b-thermal]{4b: 무차원 유한 온도 물리량}
  \framesubtitle{$t=k_BT/E_0$ · 저장된 대각화를 재사용 · 기존 켈빈 함수를 $T=t/k_B$로 감싼다 · 고정 기준 상태의 저온 근사.}
  \footnotesize\renewcommand{\arraystretch}{1.08}
  \begin{tabularx}{\textwidth}{@{}P{3.3cm}P{3.4cm}Y@{}}
    \toprule
    확인 항목 & 결과 & 범위 \\
    \midrule
    직접 공식 & $10^{-13}$ & 편극 강자성의 $F,U,S,C$, 모멘트 \\
    열역학 관계 & 수치 미분과 $10^{-7}$ & $U=F+tS$, $S=-\partial_tF$, $C=\partial_tU=t\partial_tS$ \\
    고온·저온 & 모드당 $C\to1$; 활성형 & 고온에서 $\langle n\rangle>S$는 경고와 \texttt{beyond\_lswt} \\
    ED 저온 대조 & $|F_{\mathrm{ED}}-F_{\mathrm{LSWT}}|<e^{-2\Delta/t}$ & 정사각 2$\times$3, 삼각 3$\times$3; 2-마그논부터 차이 \\
    Néel, $t>0$ & $U,C$가 $N=16$--64에서 수렴 & 모멘트는 NaN(2D 로그 발산)과 경고 \\
    기존 켈빈 계산 & $10^{-14}$ & NBCP 편극상, 같은 k·MAGSWT \\
    테스트·스냅샷 & 437 통과 / 같은 5 실패 & 새 16개; 스냅샷 차이 0 \\
    \bottomrule
  \end{tabularx}
  \par\medskip
  \supportingtext{2D에서 선형·이차 분산의 영모드는 $F,U,S,C$를 유한하게 두지만 $\langle n\rangle$을 로그 발산시킨다. 기존 보손 수 kernel은 유한 온도에서도 맞음을 독립 계산으로 확인했다.}
  \slidesource{\texttt{examples/lswt\_stage4b\_check.py}}{stage4b}
\end{frame}

\begin{frame}[label=stage4c-structure-factor]{4c: 구조인자는 모드별 무게로 내고 ED와 정확히 맞춘다}
  \framesubtitle{2026-09-30 · $S^{ab}(\mathbf q,\omega)=\sum_nW_n^{ab}(\mathbf q)\,\delta(\omega-\omega_n)+N\sum_{\mathbf G}I^{ab}(\mathbf G)\delta_{\mathbf q,\mathbf G}$ · 전체 위치 게이지, 사이트당(D26).}
  \footnotesize\renewcommand{\arraystretch}{1.08}
  \begin{tabularx}{\textwidth}{@{}P{3.3cm}P{3.2cm}Y@{}}
    \toprule
    확인 항목 & 결과 & 범위 \\
    \midrule
    ED 1-마그논 행렬 요소 & $10^{-15}$ & U(1) 편극: 정사각, DM, 2-사이트 벌집격자, $S=\tfrac32$, NBCP XXZ \\
    Néel 해석식 & $2\times10^{-13}$ & 가로 무게 $S\sqrt{(1-\gamma)/(1+\gamma)}$, 세로 0 \\
    합 규칙(확장 격자) & $10^{-16}$ & 비탄성 $\langle S(1+2n)\rangle$, 탄성 $\langle(S-n)^2\rangle$, 총합 $-S(S+1)=\langle n^2\rangle$ \\
    게이지·주기 & 정확 & $S(\mathbf q+\mathbf G_{\mathrm{prim}})=S(\mathbf q)$, 자기 역격자로는 주기 아님 \\
    결합 상관 에너지 & $E_{\mathrm{GS}}$와 $10^{-15}$ & 정사각, 삼각, NBCP Y($J_{\mathrm{PD}}$) \\
    상세 균형, Bragg & 정확 & $S^{ab}(\mathbf q,-\omega)=e^{-\omega/t}S^{ba}(-\mathbf q,\omega)$; Néel $I^{zz}(\pi,\pi)=m^2$ \\
    테스트·스냅샷 & 452 통과 / 같은 5 실패 & 새 15개; 스냅샷 차이 0 \\
    \bottomrule
  \end{tabularx}
  \par\medskip
  \supportingtext{범위: 한 마그논 가로 성분과 정확한 Bragg(2-마그논 세로 연속체는 후속). Bragg 벡터의 Goldstone 영모드에서는 비탄성 무게를 NaN으로 표시하고 탄성 항은 계산한다.}
  \slidesource{4c 검증 기록}{stage4c}
\end{frame}

\begin{frame}[label=stage4c-legacy]{4c: 기존 상관 코드는 규약이 달라 쓰지 않는다}
  \framesubtitle{\texttt{observables/correlations.py}는 연결·테스트가 없던 코드다. 새 구현과 같은 조건에서 비교했다.}
  \small
  \begin{tabularx}{\textwidth}{@{}P{3.2cm}Y@{}}
    \toprule
    항목 & 결과 \\
    \midrule
    정적 $S(\mathbf q)$ 비 & 기존/새 = 2--5.7(정사각), 3--14.4(삼각); 일정한 배수가 아님 \\
    진단 & 기존 = $N_s\times$(전체 위치 게이지 $\mathsf H(\mathbf k)$ 위에 $e^{-i\mathbf q\cdot\symbf{\tau}_i}$를 한 번 더 곱한 값); 진단식과 $10^{-7}$ 일치 \\
    원인 & 셀당 정규화(사이트당 아님)와 부격자 위상의 중복 적용 → 부격자 간 간섭 항이 틀림 \\
    부가 결함 & 부격자 지정 출력에서 \texttt{[mu+Ns, nu]}를 두 번 더함(\texttt{[mu+Ns, nu+Ns]} 누락) \\
    \bottomrule
  \end{tabularx}
  \par\medskip
  \supportingtext{새 구현은 ED와 기계 정밀도로 맞는다. 기존 코드는 수정하지 않고 비교 기록으로 남겼다. 유지·삭제는 사용자 결정이다.}
  \slidesource{4c 검증 기록}{stage4c}
\end{frame}

\begin{frame}[label=legacy-cleanup]{기존 코드 정리: 검증되지 않은 경로를 지우고 기능은 옮긴다}
  \framesubtitle{2026-09-30 · 삭제한 경로의 목적은 검증된 새 함수로 제공한다(D27).}
  \footnotesize\renewcommand{\arraystretch}{1.1}
  \begin{tabularx}{\textwidth}{@{}P{3.4cm}P{3.6cm}Y@{}}
    \toprule
    삭제 & 대체 & 확인 \\
    \midrule
    MAGSWT 격자 탐색(\texttt{opt\_method='MAGSWT'}) & \texttt{orbit\_energy\_landscape}: 임의 $\phi$ 격자의 $E_{\mathrm{cl}}$, $E_{\mathrm{qm}}$ & NBCP Y·V: $E_{\mathrm{cl}}$ 변동 $\le2\times10^{-17}$, 표본 최소가 선택값과 격자 오차 안 \\
    \texttt{observables/correlations.py} & \texttt{spin\_correlation}, \texttt{to\_ladder} & 실공간 상관의 Fourier 합 = $S(\mathbf k)$($10^{-13}$); 편극 강자성 $S^{+}S^{-}$ 무게 $2S$ \\
    MAGSWT 요청 & \texttt{ValueError}로 대체 경로 안내 & 테스트 \\
    NBCP 예제 & 고전 탐색 뒤 \texttt{select\_on\_manifold} & 예제 설정: Four MSL, \texttt{no\_degeneracy} \\
    \bottomrule
  \end{tabularx}
  \par\medskip
  테스트 457 통과 / 같은 5 실패. 회귀 스냅샷은 남은 144개 값 차이 0(삭제한 MAGSWT 탐색 항목 64개는 비교에서 빠졌다고 보고).
  \par\medskip
  \supportingtext{유지: \texttt{solve\_lswt}의 MAGSWT 정규화(D24). 보류: 실시간 상관과 지연 스펙트럼 함수는 LSWT 이론 문서의 응답 규약 검토 뒤 추가한다. 비교 결과는 2b·4c 검증 기록에 남아 있다.}
  \slidesource{기존 코드 정리 기록}{cleanup}
\end{frame}

\begin{frame}[label=quantum-fix]{\texttt{quantum} 경로의 결함 수정}
  \framesubtitle{2026-09-29 · 커밋 \texttt{d43d8d4} · classical·MAGSWT 경로와 snapshot 수치는 그대로다.}
  \small
  \begin{tabularx}{\textwidth}{@{}YYp{3.3cm}@{}}
    \toprule
    결함 & 수정 & 확인 \\
    \midrule
    L-BFGS-B가 DE보다 낮지 않으면 결과 필드가 정의되지 않고 각도가 자유 변수 배열로 남음
      & DE 각도와 그 점의 $E_{\mathrm{qm}}$로 먼저 채움(method \texttt{DE}, 전체 각도)
      & 강제 fallback 2건 \\
    비교 기준이 고전 에너지 $E_{\mathrm{cl}}$(DE)뿐임
      & $E_{\mathrm{cl}}+E_{\mathrm{qm}}$(DE)보다 낮을 때만 \texttt{DE+BFGS} 채택
      & 결과 $\le$ DE 총에너지 \\
    교란이 DE 배열을 제자리에서 바꿔 보고된 고전 각도가 어긋남
      (two\_msl: 0.06--0.10 rad, $E_{\mathrm{cl}}$ $+5.3\times10^{-4}$)
      & 복사본을 교란; 재시작마다 DE 점에서 출발
      & 2건 \\
    \bottomrule
  \end{tabularx}
  \par\medskip
  전체 테스트 225 통과 / 같은 5 실패(변경 전 220 / 5). 새 테스트 5건은 수정 전 모두 실패했다.
  수치 snapshot 208개 값은 차이 0이다. snapshot은 \texttt{quantum} 경로를 쓰지 않으므로 이 경로는 새 테스트로 확인했다.
  \par\medskip
  \supportingtext{기존 경로의 결함만 고쳤다. 전체 각도에서 $E_{\mathrm{cl}}+E_{\mathrm{qm}}$을 최소화하는 이 경로는 2b단계에서 삭제했다(D18, D22).
  \texttt{angle\_setting=None}이 TypeError를 내던 결함은 \texttt{8f4dd1e}에서 고쳤다(새 테스트 2건).}
  \slidesource{상태 선택 검증 기록}{selection}
\end{frame}

\begin{frame}[label=displacement-resolution]{결합 변위의 방향: 해석과 변환 규칙}
  \framesubtitle{D13 · 채택한 Fourier 규약에서 기존 코드는 $d=r_{\mathrm{source}}-r_{\mathrm{target}}$로 일관된다.}
  \begin{columns}[T,onlytextwidth]
    \begin{column}{0.54\textwidth}
      \small
      \begin{tabular}{@{}lccc@{}}
        \toprule
        후보 셀 & NN 잔차 $+d$ / $-d$ & $r_t=r_s-d$ & $r_t=r_s+d$ \\
        \midrule
        One & $10^{-16}$ / $10^{-16}$ & 6/6 & 6/6 \\
        Two & 0 / 0 & 12/12 & 12/12 \\
        Three & 0.3333333 / $10^{-16}$ & 18/18 & 9/18 \\
        Four & $10^{-16}$ / $10^{-16}$ & 24/24 & 24/24 \\
        \bottomrule
      \end{tabular}\\[0.4em]
      \supportingtext{잔차: 2026-09-23, 가장 가까운 정수 셀까지의 최대 차이.
      포함 수: 2026-09-29, NN·NNN 물리 결합을 한 번씩 덮는 수. Three MSL의 $+d$ 해석은 9개가 격자 밖이다.}
    \end{column}
    \begin{column}{0.42\textwidth}
      \begin{block}{판단}
        \texttt{SpinSystem} 설명은 source$\to$target이다. 코드 위상 $e^{-\mathrm i\mathbf k\cdot d}$가
        규약의 $e^{+\mathrm i\mathbf k\cdot\symbf\Delta}$와 같으려면 $\symbf\Delta=-d$이다.
        틀린 것은 설명이다.
      \end{block}
      \begin{block}{변환 규칙}
        $n=(r_s-d-r_t)A^{-1}$을 정수로 확인해 \texttt{cell\_offset}에 쓰고, $J$와 source·target은 유지한다.
      \end{block}
    \end{column}
  \end{columns}
  \par\smallskip
  \supportingtext{NBCP의 DM 항은 호출자가 설정하지 않아 0이므로 $J$는 대칭이다.
  2단계에서 변환 전후 $\mathsf H(\mathbf k)$가 원소 단위로 같아(최대 $1.7\times10^{-16}$) 확정되었다.}
  \slidesource{현재 코드 재검사와 geometry\_audit; 2026-09-29 Fourier 규약 결정}{contract,geometry}
\end{frame}

\begin{frame}[label=next-steps]{이행 단계별 확인 조건}
  \framesubtitle{순서: 이름 변경 $\to$ 공통 자료형 $\to$ NBCP 연결 $\to$ 상태 선택 $\to$ 벤치마크 $\to$ 물리량 $\to$ 위상량}
  \small\renewcommand{\arraystretch}{1.1}
  \begin{tabularx}{\textwidth}{@{}p{0.6cm}p{2.6cm}Y@{}}
    \toprule
    단계 & 작업 & 통과 조건 \\
    \midrule
    0 & 이름 변경(완료) & 이행 전후 테스트 결과 동일(현재 217 통과 / 5 실패), 40건·4회 수치 차이 0, 옛 import·저장 객체 복원 \\
    1 & 자료형·검증(완료) & 세 모델의 작은 fixture를 모델명 분기 없이 같은 소비자가 읽음; 공통 검증 사례 통과 \\
    2 & NBCP 연결(완료) & 1--4 MSL 수치 보존; $r_{\mathrm{target}}=r_{\mathrm{source}}-d$ 변환 후 $\mathsf H(\mathbf k)$ 원소 단위 일치; $k,-k$ 대응 \\
    2b & 상태 선택(구현) & NBCP Y·V의 축·선택이 독립 스캔 재현; D19 두 모드; MAGSWT 비교 \\
    3 & 최소 ED(구현) & 일반 ED가 독립 구현과 일치; 편극상 1-마그논 = 같은 토러스의 LSWT(모든 $\mathbf k$); 전개 규칙 공통 \\
    4 & 물리량(구현) & 4a 바닥상태, 4b 열역학·영모드, 4c 구조인자·결합 상관 \\
    5 & 위상량 & 하이젠버그 null test; 0이 아닌 기준 모델 비교; 기존 real-space volume 이슈 해소 \\
    \bottomrule
  \end{tabularx}
  \par\smallskip
  \textbf{공통 검증에 반드시 포함할 사례}\\[0.3em]
  잘못된 단위 · 중복 Zeeman 항 · 허용 밖 kind · 없는 endpoint · 역방향 중복 · 같은 사이트 반복 · 분수 셀 이동\\
  비대칭 $J$ · 셀 간 같은 basis 결합 · 배열 변조 · 작은 주기계의 항 중복도 · 비정합·퇴화 초격자의 상태
  \par\medskip
  \supportingtext{스펙트럼 일치만으로 올바른 변환이라고 판정하지 않는다.}
  \slidesource{구현 단계와 통과 조건}{contract}
\end{frame}

\begin{frame}[label=open-questions]{미결 사항}
  \framesubtitle{해당 단계에 도달하기 전에 결정한다.}
  \footnotesize\renewcommand{\arraystretch}{1.05}
  \begin{tabularx}{\textwidth}{@{}Yp{3.2cm}@{}}
    \toprule
    질문 & 결정 시점 \\
    \midrule
    키타에프 flux 섹터와 정확해 대조, ED 2-마그논 섹터 & 해당 벤치마크 시 \\
    키타에프 편극상을 0이 아닌 thermal Hall 기준 모델로 쓸지 & 5단계 \\
    TN 관련 상태·계산 요청·결과의 필드 규약 & TN 도입 시 \\
    후속 kind의 계수와 의미: $S$에 따른 환원, $S\ge1$ 사중극자 자유도의 LSWT 처리; 토러스에서 접히는 결합의 처리 & 고차 항 도입 시 \\
    D19 해상도 계수(2b 구현값)의 확정 & 사용자 검토 \\
    D28 세부(\(\Gamma\) 최소 반환, 판정 이름, 단열 비 경고값)와 구현 시점 & 사용자 결정 \\
    D17 확장: 사이트별 회전축, 여러 차원 영공간(삼각격자 $0<h<h_{\mathrm{sat}}$에 필요) & 해당 벤치마크 전 \\
    \bottomrule
  \end{tabularx}
  \par\medskip
  \supportingtext{결합 변위 해석(D13), 상태·요청·결과의 공통 구조(D14), 벤치마크 위치(D15), 영점 에너지 상태 선택(D17)과 그 구현 계획·판정 기준·2b 구현 방식(D18, D19, D22)은 2026-09-29, ED 범위·토러스 규칙(D23)은 2026-09-30 결정했다.
  변위 해석은 2단계의 $\mathsf H(\mathbf k)$ 원소 단위 비교로 확정했다.}
  \slidesource{규약 문서의 Open Questions}{contract}
\end{frame}
